\begin{figure}[H]
    \centering
    \caption{Correspondência biunívoca no Paradoxo de Galileu}
    \label{fig:paradoxo_galileu}

    \begin{tikzpicture}[
            scale=1,
            transform shape,
            font=\sffamily,
            >=stealth
        ]

        \def\ColorLink{\ifbool{darkmode}{DarkModeLink}{LightModeLink}}
        \def\ColorText{\ifbool{darkmode}{DarkModeText}{black}}

        % Rótulos das linhas
        \node[\ColorText, font=\bfseries, anchor=east] at (0, 1.5) {Naturais (\(\mathbb{N}\)):};
        \node[\ColorText, font=\bfseries, anchor=east] at (0, 0) {Quadrados perfeitos:};

        % Elementos da linha de cima (Naturais)
        \node[\ColorText] (n1) at (1.5, 1.5) {\(1\)};
        \node[\ColorText] (n2) at (3.0, 1.5) {\(2\)};
        \node[\ColorText] (n3) at (4.5, 1.5) {\(3\)};
        \node[\ColorText] (n4) at (6.0, 1.5) {\(4\)};
        \node[\ColorText] (ndots1) at (7.5, 1.5) {\(\dots\)};
        \node[\ColorText] (nn) at (9.0, 1.5) {\(n\)};
        \node[\ColorText] (ndots2) at (10.5, 1.5) {\(\dots\)};

        % Elementos da linha de baixo (Quadrados)
        \node[\ColorText] (q1) at (1.5, 0) {\(1\)};
        \node[\ColorText] (q2) at (3.0, 0) {\(4\)};
        \node[\ColorText] (q3) at (4.5, 0) {\(9\)};
        \node[\ColorText] (q4) at (6.0, 0) {\(16\)};
        \node[\ColorText] (qdots1) at (7.5, 0) {\(\dots\)};
        \node[\ColorText] (qn) at (9.0, 0) {\(n^2\)};
        \node[\ColorText] (qdots2) at (10.5, 0) {\(\dots\)};

        % Setas de mapeamento
        \draw[<->, \ColorLink, thick, dashed] (n1) -- (q1);
        \draw[<->, \ColorLink, thick, dashed] (n2) -- (q2);
        \draw[<->, \ColorLink, thick, dashed] (n3) -- (q3);
        \draw[<->, \ColorLink, thick, dashed] (n4) -- (q4);
        \draw[<->, \ColorLink, thick, dashed] (nn) -- (qn);

    \end{tikzpicture}

    \fonte{Elaborado pelo autor (2026).}
\end{figure}